\documentclass[11pt]{amsart}

\usepackage{amsfonts}
\usepackage{amssymb}
\usepackage{graphicx,color}
% \usepackage{a4wide}
\usepackage{amsmath}
\usepackage{amsthm}

\textwidth 14.8cm \textheight 19.5cm \topmargin 0in
\oddsidemargin 0.5in
\evensidemargin 0.5in
\parskip 1mm
\setlength{\parindent}{0pt}
\def\ds{\displaystyle}

% margin note
\newcommand\mnote[1]{\leavevmode\marginpar{\scriptsize{ #1 }}}
\newcommand\note[1]{\mnote{\textcolor{orange}{ #1 }}}
% margin note marked fixed
\newcommand\fixednote[1]{}
\newcommand\discuss[1]{\mnote{ #1 \textcolor{red}{discuss}}}
\newcommand\reword[1]{\textcolor{orange}{ #1 }}

\newcommand\lang{ILD }

\title[Title]{\bf A module system implemented on top of a minimal Scheme dialect}

\author[Angel]{Angel Angelov}
% \address{Baba \\
% foo
% }
\email{angel@qtrp.org}

% \date{\today}

\newtheorem{theorem}{Theorem}[section]
\newtheorem{lemma}[theorem]{Lemma}
\newtheorem{corollary}[theorem]{Corollary}
\newtheorem{proposition}[theorem]{Proposition}
\newtheorem{conjecture}[theorem]{Conjecture}
\theoremstyle{definition}
\newtheorem{defn}[theorem]{Definition}
\newtheorem{example}[theorem]{Example}
\newtheorem{remark}[theorem]{Remark}

\begin{document}

\maketitle

\begin{center}
Faculty of Mathematics and Informatics, Sofia University
\end{center}

\section*{Abstract}

We define a minimal language called \lang that syntactically resembles LISP
and Scheme. Functions and macros in \lang are indistinguishable and are
differentiated at the callsite. We show that the minimal features of \lang
are sufficient to define a module system that is as expressive as existing
module systems in Scheme.\cite{r6rs}

Well-established module systems for Scheme and Racket take extra care
to handle compile-time code execution and imports and exports of macros
\cite{module-system-for-scheme}\cite{submodules-in-racket}. In contrast,
we demonstrate that such techniques are not needed for a language like
\lang because of macros having runtime semantics and being effectively
identical to functions.

Using runtime macro semantics for complex metaprogramming constructs such
as the here-established module system incurs a significant performance
penalty. We propose using partial evaluation to implicitly shift module
resolution into compile-time execution, achieving similar performance
to established module systems in other languages.

\bigskip
\bibliography{citations}
\end{document}
